\documentclass[11pt]{article}
\usepackage[margin=2.3cm]{geometry}
\usepackage{amsmath,amssymb,amsthm}
\usepackage[T1]{fontenc}
\usepackage{enumitem}
\usepackage{booktabs}
\usepackage{graphicx}
\usepackage{hyperref}
\newcommand{\WIPHASH}{0237120}
\newcommand{\tri}{\mathbin{\triangleleft}}
\newcommand{\Poly}{\mathbf{Poly}}
\newcommand{\Cat}{\mathbf{Cat}}
\newcommand{\Set}{\mathbf{Set}}
\newcommand{\Comon}{\mathbf{Comon}}
\newcommand{\E}{\mathcal{E}}
\setlength{\parskip}{4pt}\setlength{\parindent}{0pt}
\title{Review: \emph{The Day convolution of two polynomial comonoids is the product internal
category --- in every Grothendieck topos}}
\author{Rick\\[2pt]\small To: Scot MacBeth (Kodamai) \quad cc: Robin}
\date{2026-10-23}
\begin{document}
\maketitle
\thispagestyle{empty}
\begin{center}
\resizebox{\linewidth}{!}{\begin{tabular}{ll}
Reviewed: & \texttt{2026-10-22-day-convolution-topos-uniform.pdf}, content \texttt{b553b27}, 8 pp.; email UID 440\\
This review: & corresponds to WIP \texttt{\WIPHASH} (\texttt{grandpa-rick/work-in-progress}, \texttt{reviews/})\\
Sources read: & arXiv abstracts of all five references; full text of Shapiro--Spivak 2305.00167 (SS),\\
 & Gambino--Kock 0906.4931 (GK), Niu--Spivak 2312.00990 (NS). Lean file not run.
\end{tabular}}
\end{center}

Counts: \textbf{3 MUST, 4 SHOULD, 6 NIT}.

\section*{Verdict}
\begin{itemize}[leftmargin=*]
\item \textbf{Theorem 1.1 is true, and \S3 is sound internally.} No pointwise or Set-based
  assumption is used (your question (i); details in S2). It is also true in much greater generality
  than you state, and its Set case is already in print (M3, S1).
\item \textbf{Theorem 1.2 is false, already for $\E=\Set$.} Comonoid morphisms in $\Poly_\E$
  are \emph{cofunctors}, not functors. In $\Cat^\sharp$ the product category is \emph{not} the
  categorical product. Niu--Spivak say so explicitly. (M1)
\item \textbf{The Gambino--Kock attribution is wrong.} (Your question (ii).) GK's distributive law
  is $\Pi$ over $\Sigma$. Your $(\star)$ is a finite-limit fact and needs no LCCC at all. The
  ``LCCC + extensive'' folklore you sharpen has no source. (M2)
\item \textbf{\S7 draws the line in the wrong place.} (Your question (iii).) Theorem 1.1 holds in any
  finitely complete $\E$ with exponentiable domain maps. That covers every elementary topos and every
  LCCC, with no coproducts needed. The one genuinely base-sensitive claim is the word ``Day
  convolution'' in the title. That is cited from a Set-only proposition and is unproved over $\E$. (M3)
\item \textbf{Recommendation: not ready for publishable-result.} After the fixes, what remains is a
  correct, honestly credited short note (proposition level): ``the Dirichlet tensor of internal
  categories is the product category, internally to any finite-limit category''. I would read a
  second round gladly.
\end{itemize}

\section*{Findings}
\subsection*{MUST}
\begin{enumerate}[label=\textbf{M\arabic*.},leftmargin=*]
\item \textbf{Theorem 1.2, \S4, abstract, \S2.2 and Table 1: the universal property is false.}
  SS Cor.\,5.12 (checked in the full text) is an isomorphism between $\Comon(\Poly_\E)$ and
  \emph{internal categories with exponentiable source and internal cofunctors} (SS Def.\,5.11). It is
  not an equivalence with $\Cat(\E)$ (functors). So ``the internal equivalence
  $\mathrm{cat}(\E)\simeq\Cat(\E)$'', and ``since $p_C$ is a functor, its forward image is a
  comonoid morphism'' (\S4), are both wrong. In $\mathrm{cat}(\E)$ the product must be computed among
  cofunctors, and there it fails.

  \emph{Literature.} NS \S8.2.4, opening sentences: ``The usual product of categories
  is not the categorical product in $\Cat^\sharp$. It is, however, a monoidal product on
  $\Cat^\sharp$, coinciding with the parallel product $\otimes$.'' NS's chapter summary adds that the
  $\Cat^\sharp$-product of the walking arrow with itself has infinitely many objects.

  \emph{Explicit witness ($\E=\Set$).} Let $C=D=\mathbf 2=\{0\xrightarrow{a}1\}$, and let $E$ have
  objects $e,x,z$ and non-identity arrows $u\colon e\to x$, $v\colon e\to z$. Define cofunctors
  $F,G\colon E\to\mathbf 2$ by $e\mapsto0$, $x,z\mapsto1$, lifting $a$ at $e$ to $u$ (for $F$) and
  to $v$ (for $G$). The cofunctor laws hold trivially: there are no composable non-identity pairs.
  Any $H\colon E\to\mathbf2\times\mathbf2$ with $\pi_CH=F$ and $\pi_DH=G$ needs $H(e)=(0,0)$ and
  $H(x)=H(z)=(1,1)$. It must then lift $(a,\mathrm{id})\colon(0,0)\to(1,0)$ to an arrow out of $e$
  whose codomain maps to $(1,0)$. No such arrow exists. So $(c\otimes d,\pi_C,\pi_D)$ is not a product.
  (Your $\pi_C,\pi_D$ \emph{are} cofunctors, lifting $f\mapsto(f,\mathrm{id})$. The cone exists; it
  is the universality that fails.)

  \emph{Fix.} Delete Theorem 1.2, its row in Table 1, and the abstract sentence. Replace them with
  a remark: $\otimes$ is a monoidal, non-cartesian product on $\Cat^\sharp(\E)$, generalizing NS
  Prop.\,8.79. Also correct the description of SS in \S2.2, \S4 and Rem.\,4.1: ``isomorphism with
  internal categories \emph{and cofunctors}''. If you want a universal property, the honest one is
  in $\Cat(\E)$ (functors), where $C\times D$ is the product by \S4's first paragraph. But that is a
  statement about $\Cat(\E)$, not about $\otimes$ on comonoids.

\item \textbf{\S1 bullets, \S2.1, \S3 Step 2, \S5 and Rem.\,4.1: $(\star)$ is not the Gambino--Kock
  law, and the ``sharpening'' is of an unsourced claim.} GK \S1.2, eq.\,(5)--(6), is the
  distributivity of \emph{dependent products over dependent sums}:
  $\Pi_f\Sigma_u\cong\Sigma_v\Pi_g\Delta_e$. In type theory this is the ``type-theoretic axiom of
  choice'', $\prod\sum\cong\sum\prod$. It is what composition of polynomials uses. Your
  $(\star)$, $\Sigma_{(i,k)}A_i\times B_k\cong\Sigma_iA_i\times\Sigma_kB_k$, says something else:
  the total space of a product of bundles is the product of total spaces. Concretely, at the point of
  use both sides are
  $(C_1\times_{C_0}C_1)\times(D_1\times_{D_0}D_1)\cong(C_1\times D_1)\times_{C_0\times D_0}(C_1\times D_1)$.
  That holds in \emph{any finitely complete} category: limits commute with limits. It needs no
  $\Pi$, no LCCC, and no colimit preservation. The gloss ``$A\times(-)=\Delta_A^*$ is a left adjoint
  and so preserves the relevant colimits'' is beside the point, because no colimit occurs. (Also,
  $\Delta_A^*$ lands in $\E/A$. The endofunctor $A\times(-)$ is $\Sigma_A\Delta_A$.)

  On ``LCCC + extensive'': the only extensivity in GK is \S4.3, for free monads and W-types (the
  coproduct $X+P(TX)$). It has nothing to do with $\tri$ or $\otimes$. Unless you can quote a source
  that asserts extensivity is needed for $c\otimes d=C\times D$, there is nothing to sharpen.
  \emph{Fix:} cite GK only for $\Poly(\E)$ and $\tri$ over an LCCC (that attribution is correct).
  Call $(\star)$ ``products commute with pullbacks''. Delete the sharpening narrative from the
  abstract, \S1 and \S5.

\item \textbf{Title, abstract, Def.\,2.2 and \S7: wrong scope line; ``Day convolution'' is unproved
  over $\E$.}
  \begin{itemize}[leftmargin=*]
  \item \emph{The theorem is far more general than stated.} SS work throughout in an arbitrary
    \emph{finitely complete} $\E$: polynomials are exponentiable maps, $\tri$ is defined (SS
    Thm.\,3.26), $\otimes$ is the product of maps (SS Def.\,4.1, preserved exponentiability by Lemma
    4.3), normal duoidality holds (SS Thm.\,4.5), and comonoids are internal categories (SS
    Thm.\,5.6). Your \S3 uses nothing beyond this. So Theorem 1.1 holds in every finitely complete
    $\E$ for internal categories with exponentiable $\mathrm{dom}$. In particular it holds in every
    LCCC and every elementary topos, Grothendieck or not, with or without a natural numbers object.
    \S7's ``a bare LCCC without the requisite coproducts is not covered'' is false. Dependent sums
    $\Sigma_f$ are postcomposition and exist in any category, and GK polynomials need no coproducts.
    The ``open upper bound'' is not open: it is the SS setting.
  \item \emph{The base-sensitive part is the Day-convolution identification.} NS Prop.\,3.79 is
    about $\Set^\Set$. Its proof is a co-Yoneda computation with coends in $\Set$. Def.\,2.2's
    ``by Niu--Spivak this $\otimes$ is the Day convolution on $\Poly(\E)$'' is unsupported. Over $\E$
    you would need to say Day convolution \emph{of what}: $\E$-valued functors on $\E$? $\E$-indexed or
    enriched? And you would need the coends $\int^{A,B}$ over the large category $\E$, with copowers by
    hom-\emph{sets} or internal homs, plus a pullback-stability argument. An elementary topos has no
    infinite coproducts, so this is exactly where Grothendieck vs.\ elementary could matter. Your
    proof never uses the Day description: $\otimes$ is \emph{defined} as the product of bundles.
  \item \emph{Fix.} Either (a) retitle to ``the Dirichlet tensor'', state Theorem 1.1 over a
    finitely complete $\E$ (exponentiable $\mathrm{dom}$), and keep ``Day convolution'' only as the
    Set-level motivation via NS Prop.\,3.79; or (b) prove the Day-convolution identification over
    Grothendieck $\E$ in a stated sense. Then that, and only that, is where the Grothendieck
    hypothesis lives. I recommend (a). Also rewrite \S7 and Table 1 accordingly.
  \end{itemize}
\end{enumerate}

\subsection*{SHOULD}
\begin{enumerate}[label=\textbf{S\arabic*.},leftmargin=*]
\item \textbf{Novelty and locators (Theorem 1.1, \S1, \S2.3).} The Set case is NS
  \textbf{Prop.\,8.79}: carrier $c\otimes d$, counit $\epsilon_c\otimes\epsilon_d$,
  comultiplication via the duoidal interchanger, and ``codomains and composition are defined
  coordinate-wise''. That is your Theorem 1.1 with $\E=\Set$, word for word. Cite it by number
  instead of ``[3, the Set case]''. Internally, SS do not state $c\otimes d=C\times D$, so the
  internal statement is fine to claim. But it is a short corollary of SS Thm.\,4.5 and Thm.\,5.6, and
  should be called one. SS's own technique gives a three-line alternative proof worth a remark:
  their Lemma 5.5 transports both comonoid and category structures along the fully faithful
  singular functor $\E\to\widehat{\E}$. This reduces to presheaves, and in presheaves the category
  side is pointwise. Locators: the interchanger is SS Thm.\,4.5, eq.\,(16), with the Set formula in
  Ex.\,4.6. Your ``[4, Ex.\,2.9, eq.\,(4)]'' points at the example ``$\mathrm{id}_\E$ is dense''.
\item \textbf{\S3 Step 3: internal validity (your question (i)).} I find no smuggled
  assumption. No choice is used: identities and composition are given maps, and the only ``choices''
  are chosen pullbacks. No excluded middle is used: there are no case splits and no decidability.
  No coend is computed pointwise: none is computed at all, since $\otimes$ is the product of maps.
  The element notation is legitimate as generalized elements, i.e.\ the internal language of a
  finite-limit category with exponentiable maps. But three points should be stated rather than
  glossed:
  \begin{itemize}[leftmargin=*]
  \item The forward map of $\delta_c$ is not ``$\mathrm{cod}$''. It is the
    $\Pi_{\mathrm{dom}}$-transpose of $\mathrm{cod}$, a map $C_0\to\Sigma_{C_0}\Pi_{\mathrm{dom}}(C_0\times C_1)$.
  \item On positions, $\iota$ is the canonical map
    $\Pi_X A\times\Pi_Y B\to\Pi_{X\times Y}(A\times B)$ (fibrewise $A^X\times B^Y\to(A\times B)^{X\times Y}$).
    Say in one line that it sends the pair of transposes of $\mathrm{cod}_C,\mathrm{cod}_D$ to the
    transpose of $\mathrm{cod}_C\times\mathrm{cod}_D$ (naturality of transposition).
  \item ``On the nose'' means equality in $\Poly_\E$, whose morphisms are \emph{isomorphism classes}
    of diagrams (SS Def.\,3.7, 3.13). That is what makes ``$=$'' meaningful despite the canonical
    isomorphism in the composition bullet. Rephrase ``there is no coherence isomorphism to chase''.
    Likewise, Rem.\,3.1's ``laws inherited from $C\times D$'' needs the fact that SS's proof of
    Thm.\,5.6 matches the \emph{equations} one by one on unconstrained data. Their statement alone
    is a bijection of lawful structures. Cite the proof.
  \end{itemize}
\item \textbf{\S1, Lucyshyn-Wright [2].} The abstract of 2502.18557 is about functor categories
  and bifunctors over \emph{non-symmetric} bases $V$, via graded categories. Nothing in it says
  ``Day-convolution closure is special to Set-like bases''. Quote a locator that does, or delete the
  sentence. The same goes for the Vec paragraph: ``Vec is not LCCC'' is true, but SS need only finite
  limits and exponentiable maps, so ``relocates the Vec obstruction precisely'' overstates.
\item \textbf{\S6, Graph witness.} This is a fine finite sanity check of Theorem 1.1 in a presheaf
  topos. But in a presheaf topos the category side is computed pointwise, so it cannot test the
  ``non-pointwise'' worry. The global-element point needs no test: $\mathrm{id}$ is a morphism, and
  global elements never enter. Present it as an illustration. Drop ``to the author's knowledge, the
  first Lean certificate\dots''. Give the Lean file's URL (the PDF has only a path).
\end{enumerate}

\subsection*{NIT}
\begin{enumerate}[label=\textbf{N\arabic*.},leftmargin=*]
\item Def.\,2.1: $p(X)=\Sigma_{A_p}\Pi_{s_p}A_p^*X$ is ill-typed. $\Pi_{s_p}$ needs an object over
  $B_p$, so the formula should be $\Sigma_{A_p}\Pi_{s_p}B_p^*X$ (i.e.\ $\Delta_{B_p}X=X\times B_p$).
\item \S6: $C_{\mathrm{pt}}$ is a baseline (positive) control, not a ``negative control''. The negative
  controls are \texttt{bogus} and \texttt{notcat}.
\item [5, Prop.\,5.2] for the comonoid phrasing: I did not verify this locator. SS Thm.\,5.6 with
  Rem.\,5.7 is a safer citation.
\item \S1 says ``Theorem 1.1 --- ours'' while also crediting ``[3, the Set case]''. Make it
  consistent once S1 is done.
\item Table 1, row 1: the column ``holds in'' should say ``any finitely complete $\E$ (exponentiable
  maps)''. See M3.
\item SS eq.\,(16) itself has a typo ($p_2\tri q_2$ for $p_2\otimes q_2$). If you reproduce it,
  don't copy that.
\end{enumerate}

\section*{What I checked, plainly}
\begin{itemize}[leftmargin=*]
\item \textbf{SS (2305.00167), full text.} Abstract: finite-limit $\E$, $\otimes$, duoidality,
  comonoids $=$ internal categories with exponentiable source. Thm.\,5.6 is the bijection of
  structures. Cor.\,5.12 is the isomorphism with \emph{cofunctors}. Your locators 5.6 and 5.12 are
  the right numbers, but the content is misdescribed (M1). Ex.\,2.9 is the wrong locator (S1).
\item \textbf{GK (0906.4931).} Polynomial functors over LCCCs, the framed bicategory, free
  monads. The distributive law is eq.\,(5)--(6), $\Pi$ over $\Sigma$. Extensivity appears only in
  \S4.3 (M2).
\item \textbf{NS (2312.00990).} Prop.\,3.79 is Day convolution on $\Set^\Set$ only. \S8.2.4 has
  Prop.\,8.79 $=$ the Set case of Theorem 1.1, and the explicit statement that $\times$ is not the
  $\Cat^\sharp$-product.
\item \textbf{Lucyshyn-Wright (2502.18557)} and \textbf{SGF (2111.10968)}: abstracts only. Titles and
  authors match your bibliography.
\end{itemize}
\end{document}
